% =====================================================================
%  GMN2 §3.2  Witten 的构造（§3.2.1 IIA 型 / §3.2.2 提升到 M 理论 / §3.2.3 共形箭图）
%  原文：wkb-paper.tex 2191–2793 行；印刷页 32–41；编号公式 (3.41)–(3.72) 共 32 条
%  本块公式编号全部落在块内，块内引用用 \eqref；跨块引用写原文编号并加 % TODO
%  图：witten-3 = 图 5、witten-2 = 图 6、witten-1 = 图 7（v0.1 硬编码，合稿改 \ref）
% =====================================================================

\subsection{Witten 的构造}
\label{subsec:WittenConstruction}

我们在第 3.1 节% TODO 跨块引用
讨论过的那套构造有一个重要的特例，多年前就由 Witten 考虑过\cite{Witten:1997sc}。确实，设想我们考虑一个 $\CN=2$、$d=4$ 理论，它由某个共形或渐近自由的酉群线性箭图（也许还带基本表示物质）所定义。在 \cite{Witten:1997sc} 中，这类箭图是用 IIA 型弦论中某些 D4/D6/NS5-膜位形来实现的。此外 Witten 还注意到，提升到 M 理论之后，这些膜位形可以自然地用单个 M5-膜来描述。

本节较为详细地回顾这一构造，并论证它可以看作第 3.1 节% TODO 跨块引用
一般图景的一个例子。特别地，在约化到 $d=3$ 之后，所得理论的模空间 $\CM$ 应当是一个 Hitchin 系统。于是我们给出一条规则：把 $\IC\IP^1$ 上某个特定的带分歧 Hitchin 系统配给任意一个共形或渐近自由的酉群线性箭图。这个 Hitchin 系统并不是最一般的可能——我们将看到，其中出现的缺陷算符属于一种受限的类型。

实际上，通过对箭图规范理论取适当进一步的标度/退耦极限，人们可以得到与 $\IC\IP^1$ 上乃至更高亏格曲线 $C$ 上更一般的 Hitchin 系统相联系的 $\CN=2$ 场论的纯四维实现。与一般黎曼面上具有正则奇性的 Hitchin 系统相联系的 $\CN=2$ 场论，其四维构造最近出现在 \cite{Gaiotto:2009we} 中。

\subsubsection{IIA 型}

沿用 Witten 的做法，我们从 IIA 型弦论出发。考虑 $n+1$ 个平行的 NS5-膜（$n\geq 1$），以 $\alpha = 0, 1, \dots, n$ 标记。这些膜沿方向 $x^{1,2,3,4,5}$ 延伸，在 $x^{7,8,9}$ 上取共同的值，并在 $x^6$ 方向上彼此分开。我们把区间
\begin{equation}
\label{eq:IIA-intervals}
I_\alpha := \{ x^6: x^6_{\alpha-1} < x^6 < x^6_{\alpha} \}
\end{equation}
称为第 $\alpha$ 个区间。接下来引入一族 D4-膜，它们位于 $x^{4,5,7,8,9}$ 的固定取值处，并沿 $x^6$ 方向延伸过某些区间、终止于 NS5-膜上。NS5-膜与 D4-膜全都位于 $x^{7,8,9}$ 的同一取值处。左侧（$x^6 \to -\infty$）有 $k_0$ 个半无穷 D4-膜，区间 $I_\alpha$ 中有 $k_\alpha$ 个 D4-膜（$1 \leq \alpha\leq n$），右侧（$x^6 \to +\infty$）有 $k_{n+1}$ 个半无穷 D4-膜。最后我们还可以引入 D6-膜，它们沿 $x^{1,2,3,7,8,9}$ 延伸，并位于 $x^{4,5,6}$ 的固定取值处。图 5% TODO 图号待合稿
画出了这样一个位形的例子。

\insfig{4.7in}{witten-3}{一个 IIA 型位形：NS5-膜（蓝）、D4-膜（紫）与 D6-膜（红圈内画叉）。这里取 $n = 2$ 与 $(k_0, k_1, k_2, k_3) = (2, 4, 3, 1)$。D6-膜位于 $x^{4,5,6}$ 的确定取值处。}

这个系统的基态保持 $8$ 个实超荷。在远大于弦长、也远大于 NS5-膜之间距离的尺度上，膜的涨落由时空 $x^\mu = x^{0,1,2,3}$ 中的一个 $d=4$、$\CN=2$ 规范理论描述。这个规范理论是一个线性箭图，规范群为 $U(k_1) \times\cdots \times U(k_n)$；见图 6% TODO 图号待合稿
。

\insfig{3in}{witten-2}{与图 5% TODO 图号待合稿
所示的 IIA 型位形相对应的线性箭图。圆形节点对应 $U(k_\alpha)$ 规范群。圆形节点之间的连线对应双基本表示物质。圆形节点与方形节点之间的连线对应基本表示物质，其重数由方形节点中的数 $d_\alpha$ 给出。}

这个箭图理论的物质内容，是每个 $U(k_\alpha)$ 因子有 $d_\alpha$ 个按基本表示变换的物质场。这些基本表示物质场有两个来源。第一，它们可以来自第 $\alpha$ 个区间内 D4-膜与 D6-膜之间伸展的弦，这时它们带 $U(k_\alpha)$ 的荷。第二，在 $x^6$ 区间任一端的那组半无穷 D4-膜与其相邻的那组 D4-膜之间伸展的弦，给出 $U(k_1)$ 的 $k_0$ 个基本表示与 $U(k_n)$ 的 $k_{n+1}$ 个基本表示。当 $n\geq 2$ 时，D4-D4 弦还给出 $U(k_\alpha) \times U(k_{\alpha+1})$ 的双基本表示超多重态，$\alpha = 1, \dots, n-1$。

正如 \cite{Witten:1997sc} 中所解释的，D4-膜对 NS5-膜施加的力与有效四维理论中各规范耦合的 β 函数之间，有一个漂亮的关系。这一关系表明，为避免大的反作用效应，我们必须把注意力限制在渐近自由或共形的理论上。因此我们要求 β 函数系数非负：
\begin{equation}
\label{eq:NegBetas}
b_\alpha:=  - 2k_\alpha  + k_{\alpha-1} +  k_{\alpha +1} +  d_\alpha \leq 0 , \qquad \alpha=1,\dots, n.
\end{equation}

规范群的中心 $U(1)_\alpha \subset U(k_\alpha)$ 不参与四维动力学。相对的 $U(1)$ 因子被希格斯化，因此商掉一个有限群之后，有效规范群是 $\prod_{\alpha} SU(k_\alpha)\times U(1)_d$；其中对角的 $U(1)_d \subset \prod_{\alpha} U(1)_{\alpha}$ 代表整族 D4-膜的总体质心坐标。这个对角的 $U(1)_d$ 是自由的，并退耦。

现在来总结这个理论的参数。UV 参数是规范耦合或 UV 标度 $\tau_\alpha$ 或 $\Lambda_\alpha$，再加上双基本表示质量 $\mu_\alpha$ 与基本表示质量 $\hat \mu_{\alpha,j}$，$1 \leq j \leq d_\alpha$（它们进入超势）。此外还有指定库仑分支上一点的 IR「$u$-参数」：记 $\Phi^{(\alpha)}$ 为 $SU(k_\alpha)$ 的矢量多重态中的伴随标量场，则对每个 $\alpha$，这些参数就是特征多项式 $\langle \det(y - \Phi^{(\alpha)}) \rangle$ 的系数。

最后我们指出，这个理论还有重要的味对称。每个双基本表示超多重态有一个 $U(1)$ 味对称（当双基本表示属于规范群 $SU(2) \times SU(2)$ 时增强为 $Sp(1)\sim SU(2)$）；每一组 $d_\alpha$ 个基本表示味有一个 $U(d_\alpha)$ 味对称（当基本表示属于规范群 $SU(2)$ 时增强为 $SO(2 d_\alpha)$）。这些味对称并不都在膜构造中显现。其中一些——例如 $SU(2)$ 规范群情形的增强——在六维 $(2,0)$ 理论框架中、把多余的 $U(1)$ 规范群因子剔除之后就会显现。另一些则只在流向四维规范理论之后才浮现。本文通篇我们将看到各种例子。

\subsubsection{提升到 M 理论}

Witten 构造的第二个关键要素，是把上面的 IIA 型位形提升到 M 理论。记 $x^{10}$ 为一条周期坐标为 $2\pi$ 的周期坐标，它参数化 M 理论圆。在没有 D6-膜时，我们取十一维 M 理论度规为
\begin{equation}\label{eq:MthryMetric}
ds^2 = \sum_{\mu=0,1,2,3} dx^\mu dx_\mu + \sum_{i=4,5,7,8,9}
(dx^i)^2 + R_{11}^2[(dx^6)^2 + (dx^{10})^2].
\end{equation}
（注意 $x^6$ 与 $x^{10}$ 是无量纲的。）有 D6-膜时，我们将把 \eqref{eq:MthryMetric} 中的某一项换成一个多中心 Taub-NUT 流形 $Q$：让 $x^{10}$ 圆非平凡地纤维化到由 $x^{4,5,6}$ 参数化的 $\IR^3$ 上。

NS5-膜与 D4-膜都提升为 M5-膜；在理论的库仑分支上，它们一般合并成单个光滑的 M5-膜，后者缠绕 $Q$ 中一条非紧的全纯曲线 $\Sigma$。对 $\Sigma$ 的约束来自一个要求：沿 $x^{10}$ 约化后必须还原出 IIA 型的图像。由此得到 $\Sigma$ 渐近形状的规则，这些规则在 \cite{Witten:1997sc} 中作了说明，并将在下一小节中回顾。理论的 IR 参数概括在定义曲线 $\Sigma$ 的多项式方程的系数之中。从 UV 参数 $\tau_\alpha,\Lambda_\alpha, \mu_\alpha, \hat \mu_{\alpha,j}, \langle \det(y^i - \Phi^{(\alpha)} ) \rangle$ 到这些 IR 参数之间有一个相当精细的映射，下一节将部分地说明它。

我们需要取 M 理论系统的一个极限，使 M 理论圆很大，即 $R_{11}/\ell \to \infty$，其中 $\ell$ 是 M 理论的长度标度。于是 IIA 型的弦耦合 $g_s = (R_{11}/\ell)^{3/2}$ 发散。尽管如此，我们仍想利用 D-膜的低能规范动力学。这是可行的，因为正如 \cite{Witten:1997sc} 第 2.3 节所解释的，规范理论的边缘与相关参数在 $g_s$ 与 $x^6$ 同时作标度变换下不变。特别地，我们希望保持四维规范耦合不变。用 D4-膜的标准 DBI 作用量，我们求得四维规范耦合 $(g_{YM}^{(\alpha)})^{-2} \sim g_s^{-1} \frac{R_{11}\Delta x^6}{\ell_s} \sim \left(\frac{\ell}{R_{11}}\right)^{1/2}(x^6_{\alpha} - x^6_{\alpha-1})$，$1\leq \alpha \leq n$，只差数值因子。因此 NS5-膜的间距为 $\Delta x^6 \to \infty$。关于这一构造适用区间的详细讨论，见 \cite{Witten:1997sc}。

我们的目的是重新解释这个膜设置，使这一构造能与前面关于 $(2,0)$ 理论的一般讨论对上：确实，它等价于在柱面上考虑一个 $A_{K-1}$ $(2,0)$ 理论，在柱面的若干有限点处带有简单缺陷，并在柱面两端带有（可能颇为复杂的）边界条件。为清楚起见，我们按复杂度递增的例子来展开。

\subsubsection{共形箭图}

最简单的设置涉及 $K$ 个无穷 D4-膜与 $n+1$ 个 NS5-膜相交。也就是说，取 $k_0 = k_1 = \cdots = k_{n+1} = K$。这对应于 $n$ 个 $SU(K)$ 规范群的线性箭图，其中第一个与最后一个因子各有 $K$ 个基本表示。

提升到 M 理论是直截了当的：$K$ 个 D4-膜提升为 $K$ 个 M5-膜，缠绕由 $s := x^6 + i x^{10}$ 参数化的柱面。这个柱面应当等同于第 3.1 节% TODO 跨块引用
的黎曼面 $C$。这 $K$ 个 M5-膜将在 $C$ 上给出一个 $A_{K-1}$ $(2,0)$ 理论。NS5-膜也提升为 M5-膜，它们在互不相同的点 $s = s_\alpha$（$\alpha=0, \dots, n$）处与柱面相交。它们的世界体充满 $x^4$–$x^5$ 平面。它们在 $C$ 上的 $A_{K-1}$ 理论中给出简单缺陷。

\insfig{3in}{witten-1}{$K$ 个 M5-膜缠绕柱面 $C$。另有 $n+1$ 个横向的单重缠绕五膜在点 $s_\alpha$ 处与 $C$ 相交。}

我们引入无量纲坐标 $t := e^{-s} \in \IC^\times$ 与 $v := (x^4 + i x^5)/\ell \in \IC$。\footnote{$t$ 与 $v$ 之间的代数关系，连同 $t$ 与 $s$ 之间的指数关系，最终导致耦合的对数式 RG 流。}图 7% TODO 图号待合稿
提示：M5-膜所缠绕的曲线 $\Sigma$ 就是 $\IC^\times \times \IC$ 中的轨迹
\begin{equation}\label{eq:degencurve}
v^K \prod_{\alpha = 0 }^n (t-t_\alpha) = 0.
\end{equation}
正如我们在描述四维规范理论解释时将会看到的，\eqref{eq:degencurve} 实际上对应共形点，此时所有质量以及所有真空期望值 $\langle {\Tr} (\Phi^{\alpha})^s \rangle $（$1\leq \alpha \leq n$，$s=1,\dots, k_\alpha$）均为零。在库仑分支上，这一设置的各个 M5-膜合并成单个光滑的黎曼面，它由 $(t,v)$ 中的一个多项式方程定义，该方程用 $v$ 的低阶项把 \eqref{eq:degencurve} 形变：
\begin{equation}\label{eq:swcurve1}
F(t,v) = v^K \prod_{\alpha =0 }^n (t-t_\alpha) + \sum_{i=1}^K p_i(t) v^{K-i} = 0.
\end{equation}
若 $p_1(t_\alpha)\not=0$，则在 $t = t_\alpha$ 处，根 $v(t)$ 当中恰有一个具有极点，留数为 $\rho_\alpha:=-\frac{p_1(t_\alpha)}{\prod_\beta (t_\alpha-t_\beta)}$。在 $t=t_\alpha$ 处对应这个发散根的 $\Sigma$ 的那一分支，在物理上被认定为 IIA 型图像中第 $\alpha$ 个 NS5-膜的 M 理论提升。

我们还希望有 $K$ 个五膜，在 $t$ 极大与极小时都趋向 $v$ 的常数值。当 $t$ 很小时，\eqref{eq:swcurve1} 的 $K$ 个根代表那 $K$ 个延伸到 $x^6 \to +\infty$ 的 D4-膜的渐近坐标。如果我们还希望在 $t \to \infty$ 的极限下 $v(t)$ 也有 $K$ 个互不相同的根，那么每个 $p_i$ 的次数至多为 $n+1$。因此 $F$ 也可以写成
\begin{equation} \label{eqn:sw1}
F(t,v)= \sum_{\alpha=0}^{n+1} q_\alpha(v) t^{n+1-\alpha} = 0.
\end{equation}
其中 $q_\alpha(v)$ 是 $v$ 的多项式，次数 $\leq K$。此外，$q_0(v)$ 与 $q_{n+1}(v)$ 的次数必须是 $K$。$q_0$ 与 $q_{n+1}$ 的根，分别是当 $t$ 极大或极小时根 $v(t)$ 的 $k_0=K$ 个与 $k_{n+1}=K$ 个渐近值。我们分别把它们记为 $v^{(0)}_1, \dots, v^{(0)}_K$ 与 $v^{(n+1)}_1, \dots, v^{(n+1)}_K$。

\subsubsection*{弱耦合}

我们希望对 $F(t,v)$ 各系数的物理含义有一些直觉。如果四维规范理论是弱耦合的，这就容易了。四维 $SU(K)$ 理论的规范耦合由两部分决定：伸缩子（它决定 D4-膜上五维规范理论的耦合），以及 NS5-膜在 $x^6$ 方向上的间距。提升到 M 理论后，这个组合只依赖于缺陷的坐标 $s$，而与柱面的总体标度 $R_{11}$ 无关。包含 θ 角的关系是
\begin{equation}
\label{eq:tau-s}
- i \pi \tau_\alpha = - i\pi \left(\frac{\theta_\alpha}{2 \pi} + \frac{4 \pi
i}{g_\alpha^2}\right) = s_{\alpha} - s_{\alpha-1}, \qquad\qquad \alpha = 1, \dots, n.
\end{equation}
这一关系基于一个假设：要读出规范耦合，先应从带缺陷的柱面上的 $(2,0)$ 理论下降到 D4-膜上的 5d SYM 世界体理论，再下降到四维低能理论，这样做是有意义的。这只在缺陷间距远大于柱面半径的极限下（即规范耦合很弱时）才成立。另一方面，在强耦合下 $\tau_\alpha$ 其实并没有精确的定义：即使理论是有限的，仍然存在依赖于重整化方案的有限重整化。最简单的选择，就是干脆把 $t_\alpha$（在总体重标度意义下）当作 $\CN=2$ SCFT 的边缘形变空间的一种方便参数化。

弱耦合下我们可以取
\begin{equation}\label{eq:wkcpl}
\frac{\vert t_{\alpha}\vert }{\vert t_{\alpha-1} \vert}  = \epsilon_\alpha \to
0, \qquad\qquad 1\leq \alpha \leq n,
\end{equation}
于是在 IIA 型的语言中 NS5-膜的间距很大。其次，令 $p_0(t) =\prod_{\alpha=0}^n (t-t_\alpha) = \sum_{s=0}^{n+1} c_{n+1-s}t^s $。那么在这个极限下 $\vert c_j\vert = \vert t_0 \cdots t_{j-1}\vert (1+ \CO(\epsilon_{j}) ) $，$j=1,\dots, n+1$。我们再令
\begin{equation}\label{eq:wkcpl2}
q_\alpha(v) = c_\alpha \tilde q_\alpha(v)
\end{equation}
断言是：首一多项式 $\tilde q_\alpha(v)$ 的根，参数化了第 $\alpha$ 组 D4-膜段的位置，也就是箭图中第 $\alpha$ 个规范群的库仑分支。确实，我们可以取 $\vert t\vert $ 落在区间
\begin{equation}
\label{eq:wkcpl-range}
\vert t_{\beta} \vert \ll \vert t \vert \ll \vert t_{\beta-1} \vert
\end{equation}
之中，使得在 IIA 型的语言里 D4-膜的 $x^6$ 位置深处于第 $\beta$ 个与第 $\beta-1$ 个 NS5-膜之间的区间内。于是
\begin{equation}
\label{eq:wkcpl-ineq1}
\vert t^{n+1} \vert \ll \vert t^{n} t_0 \vert \ll \cdots \ll \vert
t^{n+1-\beta} t_0 \cdots t_{\beta-1} \vert
\end{equation}
\begin{equation}
\label{eq:wkcpl-ineq2}
\vert t^{n+1-\beta} t_0 \cdots t_{\beta-1} \vert \gg \vert t^{n-\beta
}t_0\cdots t_{\beta} \vert \gg \cdots \gg \vert t_0 \cdots t_{n+1}
\vert
\end{equation}
从而在方程 $F(t,v) = 0$ 中主导项是
\begin{equation}
\label{eq:wkcpl-dom}
t^{n+1-\beta} t_0 \cdots t_{\beta-1} \tilde q_\beta(v) =0
\end{equation}
而 M5-膜或 D4-膜的横向位置接近于 $\tilde q_\beta(v)$ 的根。同样，在强耦合下这一近似陈述不再成立。与此同时，场论中依赖于方案的有限重整化与算符混合，使得谈论「第 $\alpha$ 个规范群的库仑分支」、并试图指认场论中特定算符的真空期望值 $\langle {\Tr} (\Phi^{\alpha})^s \rangle$，多少有点失去意义。

\subsubsection*{Seiberg–Witten 曲线}

\cite{Witten:1997sc} 的一个漂亮洞见是：曲线 \eqref{eq:swcurve1} 应当等同于四维线性箭图规范理论的 Seiberg–Witten 曲线。从我们目前的视角看，这一点可由第 3.1 节% TODO 跨块引用
的一般讨论推出。$F$ 中除 $t_\alpha$ 以外的参数，即多项式 $p_i(t)$ 的系数，编码了伴随标量的真空期望值与质量参数。

质量参数的精确指认是微妙而重要的。做到这一指认的关键，是一个观察：Seiberg–Witten 微分的留数编码了质量参数。对于手头的问题，这个微分有一个简单的标准形式，它在 \cite{Witten:1997sc} 中只是隐含的，后来在 \cite{Fayyazuddin:1997by,Henningson:1997hy,Mikhailov:1997jv} 中被明确写出。

这一计算是第 3.1.5 节% TODO 跨块引用
一般讨论的特例，但在这里更为直接，所以我们把它重复一遍。于是考虑 M5-膜的一个低能位形，其到 $\IR^{1,3}\times \IC^\times \times \IC$ 的嵌入由 $(x^\mu, t, v(t;\xi_i(x^\mu) ))$ 给出。这里 $\xi_i$ 是曲线的一组独立参数，比如说多项式 $F(t,v)$ 的系数。\footnote{注意 $t = e^{-s}$ \emph{不是}时间坐标；它是沿 $C = \IC^\times$ 的坐标。}模场 $\xi_i(x^\mu)$ 被假定为 $x^\mu$ 的缓变函数。M5-膜的作用量中有一项来自法丛标量经由诱导世界体度规产生的涨落。用 \eqref{eq:MthryMetric} 的度规 $ds^2 = dx^\mu dx_\mu + \ell^2 \vert dv\vert^2 + R_{11}^2 \vert \frac{dt}{t} \vert^2 + dx^2_{7,8,9}$，在减去参考位形中 M5-膜张力对应的能量之后，我们求得动能
\begin{equation}\label{eq:KEreduxii}
\frac{R_{11}^2}{\ell^2} \int_{\IR^{1,3}} dx^{0123}\int_{\Sigma}
\frac{dt d \bar t}{\vert t\vert^2} \frac{\partial v }{\partial \xi_i} \partial_\mu
\xi_i \overline{ \frac{\partial v }{\partial \xi_j} \partial_\mu \xi_j }.
\end{equation}
现在，由于 $\frac{\partial v }{\partial \xi_i}$ 是 $t$ 的全纯函数，我们可以这样来作 $\Sigma$ 上的积分。首先注意，如果一形式 $\frac{dt}{t}\frac{\partial v }{\partial \xi_i}$ 在 $t$ 的任一取值处（包括 $t=0, \infty$）有单极点，则积分发散。因此，$F(t,v)$ 各参数的\emph{可归一化}变分，正是那些不改变 $\lambda = v \frac{dt}{t}$ 之极点的变分。我们用 $u_i$ 标记这些可归一化参数。现在，令 $\bar \Sigma$ 为把 $\Sigma$ 嵌入射影空间并填上穿刺点后所得的紧黎曼面。这些穿刺点对应于 $t_\alpha$ 处的非紧分支，以及 $t=0,\infty$ 上方的 $K$ 个分支。引入 $H_1(\bar \Sigma;\IZ)$ 的一组基 $\gamma_a$，并令 $\CI^{ab}$ 为这组基下的（对偶）相交形式。于是我们可以算出
\begin{equation}
\label{eq:periods-intersection}
\int_{\Sigma} \frac{dt d \bar t}{\vert t\vert^2} \frac{\partial v }{\partial
\xi_i}   \overline{ \frac{\partial v }{\partial \xi_j}  }= \int_{\bar\Sigma}
\frac{dt d \bar t}{\vert t\vert^2} \frac{\partial v }{\partial \xi_i} \overline{
\frac{\partial v }{\partial \xi_j}  }=\sum_{a,b} \CI^{ab} \oint_{\gamma_a}
\frac{\partial   }{\partial \xi_i}\lambda \oint_{\gamma_b} \overline{\frac{\partial
}{\partial \xi_j}\lambda }.
\end{equation}
由此我们得到法丛标量动能项的低能有效作用量：
\begin{equation}\label{eq:dualinvtke}
\frac{R_{11}^2}{\ell^2} \int_{\IR^{1,3}} dx^{0123} \CI^{ab} \partial_\mu
\Pi_a \partial_\mu \bar \Pi_b
\end{equation}
其中 $\Pi_a = \oint_{\gamma_a} v \frac{dt}{t}  $。如果我们选取一个对偶标架 $(\gamma_I, \gamma^I)$，则 \eqref{eq:dualinvtke} 变成在 $\CN=2$ 有效作用量中常见的标准表达式 $\int {\rm Im}(\tau_{IJ}) \partial_\mu a^I \partial_\mu \bar a^J$。于是我们断定：Seiberg–Witten 微分是
\begin{equation}
\label{eq:SWdifferential}
\lambda_{SW} = \lambda = v \frac{dt}{t}.
\end{equation}

众所周知，Seiberg–Witten 微分的留数仿射线性地依赖于 UV 拉氏量的质量参数。微分 $v \frac{dt}{t}$ 在覆盖 $t=0$ 的 $K$ 个分支上与覆盖 $t=\infty$ 的 $K$ 个分支上有单极点。这些极点处的留数，分别是箭图左右两端半无穷 D4-膜在 $v$-平面中的渐近位置。我们可以如下验证它们仿射线性地依赖于质量：考虑 IIA 型图像，并设 $k_1$ 个 D4-膜位于第一个区间 $I_1$ 中共同的 $v$-位置 $v_0$。经典弦论表明，$SU(k_1)$ 基本表示的 $K$ 个场的质量为 $\ell_s^{-1}\vert v_i^{(0)} - v_0\vert $，其中 $v_i^{(0)}$ 是 $q_0(v)$ 的 $K$ 个根。由全纯性我们看到，复质量参数 $\hat \mu^{(1)}_i$ 满足 $ v_i^{(0)} - v_j^{(0)} \propto \hat \mu_{1,i} - \hat \mu_{1,j} $。对于延伸到 $t\to +\infty$ 的那些分支所对应的基本表示场，有类似的陈述。于是，每一端 $SU(K)$ 味对称所对应的质量参数之差，由 $q_0(v)$ 的根之间、以及 $q_{n+1}(v)$ 的根之间的差给出。

最后，对多项式的一般参数，正如我们已经看到的，当 $t\to t_\alpha$ 时恰有一个根 $v(t)$ 有单极点，导致 Seiberg–Witten 微分有留数为 $\rho_\alpha/t_\alpha$ 的单极点。在 \eqref{eq:wkcpl} 及其后所述的弱耦合区间里，不难证明：若写
\begin{equation}\label{eq:Paramqalpha}
q_{\alpha}(v) = c_\alpha (v^K- \mu_\alpha v^{K-1} - u_2^{(\alpha)} v^{K-2} - \cdots
- u_K^{(\alpha)} )
\end{equation}
则 $t_\alpha$ 处的留数约为 $(\mu_\alpha - \mu_{\alpha+1})$（$0\leq \alpha \leq n$），并且它们要么与箭图第 $\alpha$ 个与第 $\alpha+1$ 个节点之间双基本表示物质场的质量参数一致，要么（当 $\alpha=0$ 或 $\alpha=n$ 时）与箭图任一端基本表示物质场的总体 $U(1)$ 质量参数一致。

在场论极限下，M5-膜系统的总体质心自由度退耦。我们的目标是得到一个只涉及 $A_{K-1}$ $(2,0)$ 理论自由度的描述。忽略总体质心位移的一个自然办法，是消去 $F(t,s)$ 中 $v^{K-1}$ 的系数。% 原文如此，疑为 F(t,v) 之误
为此定义
\begin{equation}
\label{eq:tilde-v}
\tilde v := v - p(t)
\end{equation}
并作选择 $p(t) = -\frac{p_1(t)}{K\prod_\alpha (t-t_\alpha)}$。注意 $t=t_\alpha$ 处的留数是 $\rho_\alpha/K$。

其次回想，所有物理量都依赖于 Seiberg–Witten 微分沿 Seiberg–Witten 曲线中闭路的积分。在给定的 IR 理论中，Seiberg–Witten 微分 $\lambda$ 并不唯一：只要求它（在任一对偶标架中）满足
\begin{equation}
\label{eq:tau-from-period}
\tau_{IJ} = \frac{\partial}{\partial a^J} \oint_{\gamma_I} \lambda,
\end{equation}
但它可以被改动，加上一个不依赖于可归一化形变的一形式。（回想：可归一化形变这一概念，是由动能的有限性 \eqref{eq:KEreduxii} 定义的。）特别地，由于 $p(t)$ 是 $t$ 的单值有理函数，我们可以自由地用平移 $v \to v + p(t)$ 来改变 Seiberg–Witten 微分的定义。$p(t) \frac{dt}{t}$ 的留数是理论质量参数的线性组合，而一阶导数 $\lambda_{u_i}$ 不变，仍与 Seiberg–Witten 曲线上的全纯 1-形式一致。$\lambda$ 的周期则被平移了质量参数的某个线性组合。为了把中心荷写成「规范荷」与「味荷」两部分贡献之和，我们需要选取序列 (3.14)% TODO 跨块引用
的一个分裂。一旦作出这样的分裂，$\lambda$ 的平移就相当于把 BPS 粒子的味荷按其规范荷的整数倍平移，也就是对味流作一次合法的重新定义。味荷的平移并非只是无害，它实际上是一种有用的改进。例如在 $K=2$ 的情形我们看到，在场论极限下味对称群从 $U(1)$ 增强为 $SU(2)$。$\lambda$ 的平移给出 BPS 粒子的电荷赋值，它与按新的 $SU(2)$ 的不可约表示来组织相容（更多信息见第 3.2.9 节% TODO 跨块引用
）。

新的 Seiberg–Witten 微分 $\tilde v \frac{dt}{t}$ 在 $t=t_\alpha$ 上方的不同叶上有单极点，其留数列表为
\begin{equation}
\label{eq:residues-malpha}
\left( - (K-1) m_\alpha, m_\alpha, \dots, m_\alpha \right)
\end{equation}
其中
\begin{equation}\label{eq:def-m-alpha}
m_\alpha = \frac{1}{K t_\alpha} \frac{p_1(t_\alpha)}{\prod_{\beta \not=\alpha}
(t_\alpha -t_\beta)}= - \frac{\rho_\alpha}{K t_\alpha}
\end{equation}
被认定为某个 $U(1)$ 味对称的质量参数。

在 $t=0$ 处（$t=\infty$ 处类似），不同叶上的极点将具有形状
\begin{equation}
\label{eq:residues-m0}
\left( - \sum m_{0,i}, m_{0,1}, \dots, m_{0,N-1} \right)
\end{equation}
其中 $m_{0,i}$（以及 $t=\infty$ 处相应的 $m_{\infty,i}$）与半无穷 D4-膜的相对位置有关，并被认定为任一组 $K$ 个基本表示的 $SU(K)$ 味对称的质量参数。注意若 $K=2$，两类缺陷是相同的，而 $U(1)$ 确实被提升为 $Sp(1) \sim SU(2)$。

\subsubsection*{一个例子}

我们用一例来结束本节。考虑 $n=1$，于是我们有一个简单的规范群 $SU(K)$，带 $2K$ 个基本表示。这一理论的标准 Seiberg–Witten 曲线，表面上看来与这里所考虑的曲线很不一样，尤其是在把所有质量参数都包含进来时 \cite{Argyres:1995wt}：
\begin{align}
y^2 & = P(w)^2 - (1-g^2) Q(w) \label{eq:example-sw-1}\\
P(w) & = \langle \det(w-\Phi)\rangle  = w^K - u_2 w^{K-2} -\cdots
- u_K \label{eq:example-sw-2}\\
Q(w) & =\prod_{I=1}^{2K} (w + g \mu + \mu_I) \label{eq:example-sw-3}
\end{align}
这里参数 $\mu, \mu_I$ 通过 $\mu = \frac{1}{2K} \sum_{I=1}^{2K} m_I $ 与 $\mu_I = m_I - \mu$ 决定味质量。参数 $g$ 是耦合 $\tau$ 的某个模函数（见 \cite{Argyres:1995wt} 的 (5.15) 式），在弱耦合下约化为 $g= 1 + \CO(e^{i \pi \tau})$。Seiberg–Witten 微分是
\begin{equation}
\label{eq:SWstandard}
\lambda^{\rm Standard} = \frac{w + (g-1) \mu}{2\pi i} d \log \left(
\frac{y-P}{y+P} \right).
\end{equation}

通过变量变换，我们可以把曲线化成 $F(t,v)=0$ 的形式，其中 $F(t,v)$ 形如 \eqref{eq:swcurve1}。首先引入满足
\begin{equation}\label{eq:couplrel}
\frac{4}{(\gamma + \gamma^{-1})^2} = 1-g^2.
\end{equation}
的 $\gamma$。其次令 $v= w  + (g-1) \mu$；再引入 $v_I := m_I$，便有
\begin{equation}
\label{eq:Qw-as-v}
Q(w) = \prod_{j=1}^{2K} (v-v_I).
\end{equation}
对称群 $S_{2K}$ 作用在根集 $v_I$ 上。（$S_{2K}$ 是 $U(2K)$ 味对称群的 Weyl 群，而后者被质量破缺到其嘉当子代数。）现在我们显式地把对称性破缺到 $S_K \times S_K$：选取一组 $K$ 个根 $v_i$，$i=1,\dots, K$。我们把剩下的根记为 $\tilde v_i$，$i=1, \dots, K$。现在令
\begin{equation}
\label{eq:t-definition-example}
t=  \frac{1}{2}(\gamma + \gamma^{-1}) \frac{y-P(w)}{\prod_{i=1}^K( v-v_i)}
\end{equation}
最后取 $ \tilde P(v) = P(w)$。这样我们就把曲线化成了形状
\begin{equation}\label{eq:APStwo}
\prod_{i=1}^K (v-v_i) t^2 + (\gamma + \gamma^{-1}) \tilde P(v) t +
\prod_{i=1}^K (v-\tilde v_i) = 0
\end{equation}
它是 \eqref{eqn:sw1} 的形式。注意，若把 \eqref{eq:APStwo} 写成 \eqref{eq:swcurve1} 的形式，我们会发现 $t_0 = \gamma, t_1 = \gamma^{-1}$，于是由 \eqref{eq:couplrel} 我们确认 $t_0/t_1 \sim e^{i \pi \tau}$，与关于弱耦合极限的讨论相符。还要注意，在这一表述中 $p_1(t)\not=0$，所以我们尚未按标准方式固定质心自由度。

值得指出，标准的 Seiberg–Witten 微分与我们的讨论中自然出现的那一个并不相同：
\begin{equation}
\label{eq:SWstandard-vs-Hitchin}
i \pi \lambda^{\rm Standard} - \lambda^{\rm Hitchin} = \frac{1}{2} v d
\log \left( \frac{\prod_{i=1}^K (v-v_i)}{\prod_{i=1}^K (v-\tilde
v_i)} \right)
\end{equation}
差别的留数，即 $v_i$，正是质量参数，所以这同样对应于对味荷作一次简单的重新定义。在去掉质心部分之前，$\lambda^{\rm Hitchin}$ 在 $t_0= \gamma$ 或 $t_1 = \gamma^{-1}$ 处的单个叶上有极点。要穿过所有变量变换去提取留数需要一些功夫，但结果简单得出奇：$t_0$ 处的留数是 $\sum \tilde v_i$，即一侧 $K$ 个半无穷膜所对应的 $K$ 个基本表示味的总体质量参数。$t_1$ 处的留数类似地是 $\sum v_i$。

$\lambda^{\rm Hitchin}$ 在 $t=0, \infty$ 处的留数与 $v_i$ 或 $\tilde v_i$ 一致；但在我们去掉质心之后，留数被修改为 $v_i - \frac{1}{K} \sum v_i$ 与 $\tilde v_i - \frac{1}{K} \sum \tilde v_i$，它们就是作用在每一组 $K$ 个基本表示味上的 $SU(K)$ 味子群的质量参数。

这一分裂是自然的：从 $(2,0)$ 理论的角度看，$U(2K)$ 味对称是一个偶然的红外对称性。六维理论只有 $U(1)^2 \times SU(K)^2$ 味对称，而每个因子都以某种方式联系到 $t=0, t_0,t_1,\infty$ 四个穿刺点之一。注意 $0,\infty$ 处的穿刺点与 $t_0, t_1$ 处的穿刺点性质不同。下一小节我们将看到这一差别如何在 Hitchin 系统的语境中体现出来。
